% !TeX root = ../papers/nc-p3.tex
\section{異なる対称三平面での二交換子鉛筆の復元}
\label{sec:marked-two-commutators}
\begin{theorem}\label{thm:marked-two-commutators}
前節の $q=3$, $g>3$, $r=2$ の場合、$P=Q$ とし、
$Q_s,Q_t$ が可換になるパラメータ二平面を取る。
その平面と補元 $w$ に対して交換子式の既知の右辺を
$C_0=[Q_s,Q_w]$, $C_1=[Q_t,Q_w]$ と読む。この二行列の鉛筆は恒等的に特異にはならない。
従って平面の基底を取り直して $C_0$ は可逆とできる。
$G,F$ とこの基底を固定すると、残る $Q$ の全復元は次の線型演算で決まる。
\begin{enumerate}
\item $D C_0+C_0 D=C_1$ を解く。解 $D$ は一意である。
\item $D$ の最小多項式が三次、$\operatorname{tr}(D^2 C_0)=0$ を要求し、
$[D,Z]=C_0$ の一解 $Z$ を取る。$L_0$ を行列 $L$ の跡零部分と書くと、全候補は
\[
 Q_s=h D_0,\qquad Q_t=h(D^2)_0,\qquad
 Q_w=h^{-1} Z_0+eD_0+f(D^2)_0,\qquad h\ne0 .
\]
\end{enumerate}
ここで $s$ は一般に取って $Q_s$ が巡回となる方向を用いる。
この式と元の有限窓の検査で印付き候補を復元できる。これだけで全次数の正則性は主張しない。
\end{theorem}
\begin{lemma}\label{lem:commuting-linear-symmetric-nets}
跡零対称三行列の二つの線型三列 $G_x,F_x$ がともに独立で既約とし、
$[G_x,F_x]=0$ が全 $x$ に成り立つなら、$F_x=hG_x$ となる。
\end{lemma}
\begin{proof}
$G$ の一般元には相異なる三固有値がある。さもなくば判別式恒等式から、
定理\ref{thm:rank-nine-finite-symmetry}の証明の通り、その中に二次元以上の冪零平面 $N$ があり、
$G$ は $N$ と跡直交する。$N$ が全て階数1以下なら、対称な階数1全体の線型部分空間は
次元1以下なので矛盾する。正則冪零な $X\in N$ を取りJordan基底で $X=E_{12}+E_{23}$ とする。
その基底で計量は反対角単位行列にできる。実際巡回ベクトル $v$ による
$v,Xv,X^2v$ の積対合は $k[t]/t^3$ の完全積跡であり、標準の最上係数との比の単元の平方根で基底を調整できる。
$Y\in N$ への $\operatorname{tr}(X^2Y)=\operatorname{tr}(XY)=0$ は
$Y_{31}=0$, $Y_{21}+Y_{32}=0$ を与え、自己随伴性で両成分は等しいので零である。
ゆえに $N$ は厳密上三角自己随伴の二平面全体であり、$G$ も跡直交と自己随伴性から上三角になって矛盾する。

相異なる三固有値の一般 $G_x$ を直交対角化する。$F_x$ も対角である。
対角固有値をそれぞれ $a_i,b_i$ とし、$k_{ij}=(b_i-b_j)/(a_i-a_j)$ とする。
偏極により全 $y$ で $(F_y)_{ij}=k_{ij}(G_y)_{ij}$ ($i\ne j$) となる。
三つの比が二つでも等しければ全て等しく、跡零性から $F_x$ は $G_x$ の倍である。
一般の $x$ でそうなら線型性により倍率は定数である。

従って三つとも異なる場合を除けばよい。$l_{ij}(y)=(G_y)_{ij}$ とする。
$[G_y,F_y]_{ij}=0$ から $l_{ij}$ のある一次式倍は
$(k_{hj}-k_{ih})l_{ih}l_{hj}$ に等しい（$h$ は第三の添字）。
非零な $l$ の辺グラフは既約性により連結であり、二辺だけなら第三辺のこの等式に反する。
従って全て非零で、各一次式は他の二つの積を割る。三つとも比例し、
$G$ の二次元部分は対角跡零全体になる。そこでは $F$ も対角である。
全非対角成分が非零の一元とこの対角平面を偏極すれば、全ての対角元でも
$F$ の対角差は対応する $G$ の差の $k_{ij}$ 倍である。独立な二差を任意に指定し、
三番目の差が和であることを使うと全 $k_{ij}$ は等しくなって矛盾する。
\end{proof}
\begin{proof}[鉛筆と巡回中心化環]
交換子次元2なので、$\wedge^2\langle Q\rangle$ の核の直線が可換二平面を指定する。
この平面に巡回元がある。なければ、二固有値の半単純元への可換摂動を使う
定理\ref{thm:rank-eight-finite-symmetry}の議論により全元が平方零で、共通像直線または共通核平面を持つ。
そのとき二交換子 $[Q_s,Q_w],[Q_t,Q_w]$ に共通左核または共通右核が生じる。
しかしこれらは $[G_a,F_c]$ の全非零成分を $ac$ 行列として張るため、共通核なら
いずれかの非零な跡零元が相手の既約三平面全体と交換してしまう。

一般の $Q_s$ は巡回で、その中心化環は三次元なので $I,Q_s,Q_t$ の包と一致する。
もし二交換子鉛筆が全て特異なら $\det[Q_s,Q_w]=0$ であり、
定理\ref{thm:rank-eight-finite-symmetry}で示した六元の行列式恒等式により
$Q_s,Q_w$ はある非零対称形式に関して自己随伴である。$Q_t$ も $Q_s$ の多項式であり、
既約性からその形式は可逆である。適切な $U$ 基底で三つの $Q$ を全て対称にする。
この変更は $G,F$ の各三列を線型結合するのでその対称性を保つ。
すると $[Q_b,Q_d]$ は $a,c$ で交代し、交換子式で $[G_x,F_x]=0$ となる。
補題により $g=3$ となって矛盾する。

最後に $C_0$ を可逆に取れば、跡零三行列のどの二固有値も和が零にならない。
したがってSylvester作用素 $L\mapsto LC_0+C_0 L$ は可逆である（上三角化で対角に固有値の和が並ぶ）。
$Q_t=a Q_s^2+bQ_s+cI$ と書き、独立性より $a\ne0$ である。
交換子を取ると一意な $D$ は $a Q_s+bI/2$ であり、$h=1/a$ とすれば表示の二式を得る。
巡回 $D$ に対する $[D,-]$ の像は $I,D,D^2$ の跡直交補である。
$C_0$ は跡零で $2\operatorname{tr}(DC_0)=\operatorname{tr}C_1=0$ だから、残る条件は表示の一つである。
これで $Q_w$ の式も従い、逆に三つの交換子は式から直接指定値に戻る。
\end{proof}
